% ID: MAT-ALG-PER-003
% Titolo: Medaglie e ripartizione proporzionale
% Disciplina: Matematica
% Area: algebra
% Argomento: percentuali
% Classe: Prima secondaria di secondo grado
% Difficolta: 3
% Tipo: problema_numerico
% Tag: matematica, algebra, percentuali
% Fonte: esercizi_razionali_anonimizzati
% Provenienza: materiale fornito dal docente; testo anonimizzato; soluzione redatta per Exergo

\begin{esercizio}
A una gara partecipano 2200 studenti. Il $25\%$ riceve una medaglia. Le medaglie di bronzo sono il quadruplo di quelle d'argento; queste sono il doppio di quelle d'oro. Quante sono quelle d'argento? Quale percentuale dei partecipanti riceve l'oro?
\end{esercizio}

\begin{soluzione}
Le medaglie sono 550. Ponendo $o$ il numero di ori, gli argenti sono $2o$ e i bronzi $8o$. Quindi $11o=550$, $o=50$, argenti 100. La percentuale di ori è $(50/2200)100\%=25/11\%\simeq2{,}27\%$.
\end{soluzione}
